% =====================================================================
%  Lecture 25 -- Network Evolution: a Grand Synthesis    (English edition)
% =====================================================================
\documentclass[11pt,a4paper]{article}
\providecommand\ACNroot{../..}
\input{\ACNroot/templates/lecture-en}

\ACNset{number}{25}
\ACNset{title}{Network Evolution: a Grand Synthesis}
\ACNset{shorttitle}{Network evolution}
\ACNset{titleother}{تکامل شبکه: یک جمع‌بندی بزرگ}
\ACNset{part}{Part V -- Planes and outlook (Lectures 23--26)}
\ACNset{time}{6--8 hours (3 h reading, 3--5 h exercises)}
\ACNset{level}{Advanced}
\ACNset{prereq}{Lectures 1--24}
\ACNset{gen}{5}

\begin{document}
\ACNcover

% ---------------------------------------------------------------------
\section{Learning objectives}
After studying this lecture you should be able to:
\begin{objectives}
  \item \textbf{Trace} the evolution of radio access, core architecture, signalling, security and QoS from 2G to 5G.
  \item \textbf{Map} network functions across generations (MSC/VLR, MME, AMF; HLR, HSS, UDM; GGSN, P-GW, UPF).
  \item \textbf{Decompose} capacity growth into spectrum, spectral efficiency and densification.
  \item \textbf{Quantify} long-term trends in peak rate, latency and cost per bit with simple growth models.
  \item \textbf{Identify} the design principles that recur in every generation and will shape the next one.
\end{objectives}

% ---------------------------------------------------------------------
\section{Prerequisite knowledge}
This lecture draws on the whole course: GSM, GPRS and the 2G/3G core (Lectures~1--4), LTE and the EPC (Lectures~5--13), and the 5G system,
SDN, NFV and SBA (Lectures~14--24). No new protocol is introduced; the aim is to connect what you already know into one picture.

\paragraph{Key terms} The vocabulary of this lecture with its Persian equivalents:
\begin{keyterms}
Generation & \fa{نسل} & A family of standards meeting a set of ITU requirements (IMT-2000, IMT-Advanced, IMT-2020). \\
Flat architecture & \fa{معماری تخت} & Fewer nodes on the user path; radio control in the base station. \\
All-IP & \fa{تمام‌IP} & Voice and data carried as IP packets, no circuit-switched domain. \\
CP/UP separation & \fa{جداسازی صفحهٔ کنترل و کاربر} & Independent scaling and placement of control and user plane. \\
Spectral efficiency & \fa{بازدهی طیفی} & Throughput per hertz per cell (bit/s/Hz/cell). \\
Densification & \fa{متراکم‌سازی} & More cells per unit area. \\
Area traffic capacity & \fa{ظرفیت ترافیکی ناحیه} & Total throughput per unit area (bit/s/m$^2$). \\
Cost per bit & \fa{هزینه به ازای بیت} & Network cost divided by carried traffic. \\
Softwarisation & \fa{نرم‌افزاری‌شدن} & Network functions as software on general-purpose infrastructure. \\
Backward compatibility & \fa{سازگاری با گذشته} & Interworking of a new generation with installed ones. \\
\end{keyterms}

% ---------------------------------------------------------------------
\section{Introduction and motivation}
Each generation of mobile networks was designed to fix the bottleneck of the previous one: 2G digitised voice and added security and roaming;
3G brought packet data at useful rates; 4G removed the circuit domain and flattened the network; 5G made the network programmable and
service-specific. Seen one at a time, the generations look like a sequence of unrelated protocol stacks. Seen together, they show a small set of
persistent ideas -- separation of concerns, mobility anchors, subscriber databases, tunnelling, mutual authentication -- implemented with ever
more general technology \cite{itu2015m2083,ts23501}. Recognising these ideas is what lets an engineer understand a new system quickly, and
predict where it will go.

% ---------------------------------------------------------------------
\section{Fundamental concepts}

\subsection{Five dimensions of evolution}
\begin{itemize}
  \item \textbf{Radio access}: FDMA/TDMA (GSM), CDMA (UMTS), OFDMA with MIMO (LTE), flexible-numerology OFDM with massive MIMO and mmWave (NR).
  \item \textbf{Core}: circuit switching with MSC/VLR and SS7/MAP; an overlaid packet core (SGSN/GGSN); the all-IP EPC with Diameter and GTP;
        the service-based 5GC on cloud infrastructure.
  \item \textbf{Mobility}: hierarchical location and routing areas, anchors in the GGSN, P-GW and UPF; handover decisions moving from the BSC and RNC
        into the base station.
  \item \textbf{Security}: one-way GSM authentication (A3/A8) to mutual UMTS AKA, EPS-AKA with key hierarchy, and 5G-AKA/EAP-AKA$'$ with concealed
        identities (SUCI) and home control.
  \item \textbf{QoS}: from circuit guarantees through PDP contexts and EPS bearers to per-flow QoS with standardised 5QIs and slicing.
\end{itemize}

\subsection{Flattening the path}
The user-plane path has steadily shrunk (\cref{fig:arch}). In GPRS and UMTS a packet passed through a base station, a radio controller (BSC/PCU
or RNC), an SGSN and a GGSN. LTE moved radio control into the eNB and removed the controller; the S-GW and P-GW remained. In 5G a single UPF can
serve as the whole core user plane, and it can be placed close to the radio.

\begin{acnfigure}
  \resizebox{\linewidth}{!}{\acnfig{lecture-25/arch-evolution}}
  \caption{User-plane path of packet data from 2.5G to 5G (simplified).}
  \label{fig:arch}
\end{acnfigure}

\subsection{Functions that survive}
Names change, functions persist. The visitor location and mobility function (MSC/VLR, SGSN, MME, AMF), the subscriber database (HLR/AuC, HSS,
UDM/AUSF), the mobility anchor and gateway to data networks (GGSN, P-GW, UPF with SMF) and policy control (PCRF, PCF) appear in every
generation, each time with cleaner interfaces and better separation between control and user plane.

% ---------------------------------------------------------------------
\section{Detailed technical explanation}

\subsection{From circuits to packets to services}
2G was designed around a 64\,kbit/s circuit per call. GPRS added a packet core in parallel, reusing the radio with dynamic time-slot
allocation. UMTS kept two domains (CS and PS) with a common radio network. LTE dropped the CS domain: voice became an IMS service over VoLTE,
with CS fallback and SRVCC for interworking. 5G keeps IMS for voice (VoNR) but restructures the core as services consumed over HTTP/2 -- the
network itself becomes a set of APIs.

\subsection{Signalling protocols}
SS7 and MAP connected the 2G/3G core; Diameter replaced MAP in the EPC (S6a, Gx, Rx); GTPv2-C carried session control; 5G uses REST-style
service-based interfaces with OAuth-based authorisation, and PFCP for user-plane control. Each step moved towards protocols used in the wider IT
industry, lowering development cost and enabling cloud deployment.

\subsection{Security lessons}
GSM authenticated only the subscriber, allowing false base stations; UMTS added network authentication and integrity protection; LTE separated
keys per layer and per node; 5G encrypts the permanent identity (SUCI), gives the home network the final authentication decision and protects
inter-operator signalling with the SEPP. Each fix answered an attack that became practical on the previous generation.

\subsection{Coexistence}
No generation replaces its predecessor at once. Networks run several in parallel for a decade or more, with interworking functions (N26 between
AMF and MME, EN-DC in non-standalone 5G, CS fallback) and gradual spectrum refarming. Backward compatibility is a constant design constraint and a
large part of the cost of every new system.

% ---------------------------------------------------------------------
\section{Mathematical models}

\subsection{Peak-rate growth}\label{sec:peak}
\paragraph{Model} An exponential trend $R(t) = R_0\, e^{k (t - t_0)}$ fitted through two points $(t_0, R_0)$ and $(t_1, R_1)$ has
\begin{equation}\label{eq:growth}
  k = \frac{\ln (R_1/R_0)}{t_1 - t_0}, \qquad T_2 = \frac{\ln 2}{k} .
\end{equation}
\paragraph{Example} From GSM data (9.6\,kbit/s, 1991) to the IMT-2020 requirement (20\,Gbit/s, specified 2017--2018) \cite{itu2017m2410}:
$k = \ln(2.08\times10^{6})/27 = 0.539$ per year, a factor 1.71 per year, doubling every 1.29 years (\cref{fig:peak}).
\paragraph{Limitation} Peak rates are specification values under ideal conditions; user-experienced rates grow more slowly and depend on load and
coverage.

\begin{acnfigure}
  \acnfig{lecture-25/peak-rates}
  \caption{Specified downlink peak rates of 3GPP technologies and the exponential fit of \cref{eq:growth}.}
  \label{fig:peak}
\end{acnfigure}

\subsection{Decomposing capacity}\label{sec:cap}
\paragraph{Model} The area traffic capacity is the product of cell density $n$ (cells/km$^2$), bandwidth $B$ and spectral efficiency $\eta$:
\begin{equation}\label{eq:area}
  C_A = n\, B\, \eta ,
\end{equation}
so the growth between two deployments is the product of the three ratios.
\paragraph{Example} LTE baseline: $n = 4$, $B = 20$\,MHz, $\eta = 1.5$\,bit/s/Hz: $C_A = 120$\,Mbit/s/km$^2$. 5G deployment: $n = 16$, $B = 100$\,MHz,
$\eta = 6$: $C_A = 9.6$\,Gbit/s/km$^2$ -- a factor $4\times5\times4 = 80$ (\cref{fig:cap}). For comparison, IMT-2020 requires 10\,Mbit/s/m$^2$
($10$\,Tbit/s/km$^2$) in the indoor hotspot scenario \cite{itu2017m2410}.

\begin{acnfigure}
  \acnfig{lecture-25/capacity-decomp}
  \caption{Area-capacity growth from the LTE baseline to the 5G deployment of \cref{sec:cap}.}
  \label{fig:cap}
\end{acnfigure}

\subsection{Latency along the path}\label{sec:path}
\paragraph{Model} The one-way user-plane latency through $h$ nodes is
\begin{equation}\label{eq:path}
  T = \sum_{i=1}^{h} \left( t_{\mathrm{proc},i} + \frac{d_i}{v} \right) ,
\end{equation}
with processing time $t_{\mathrm{proc},i}$, link length $d_i$ and propagation speed $v \approx 2\times10^{8}$\,m/s in fibre.
\paragraph{Example} Four nodes with 0.5\,ms each and 400\,km of fibre (2\,ms) give 4\,ms; two nodes (gNB and edge UPF) with 0.25\,ms each and 20\,km give
0.6\,ms. Flattening and edge placement together remove most of the core latency.

\subsection{Cost per bit}\label{sec:cost}
\paragraph{Model} If traffic grows by a factor $1+g$ per year and revenue by $1+r$, the cost per bit must fall by $(1+r)/(1+g)$ per year to keep margins
constant; it halves every
\begin{equation}\label{eq:cost}
  T_{1/2} = \frac{\ln 2}{\ln (1+g) - \ln (1+r)} \ \text{years}.
\end{equation}
\paragraph{Example} With $g = 30\,\%$ and flat revenue, $T_{1/2} = 2.64$ years. This pressure, more than peak rate, explains the industry's move to
commodity hardware, automation and software-defined networks.

% ---------------------------------------------------------------------
\section{Network architecture and diagrams}
\Cref{fig:arch} shows the user path; the control side evolved in parallel from MSC/VLR and SGSN to the MME and then to the AMF and SMF of the
service-based core. \Cref{tab:map} maps the functions across generations.

\begin{table}[htbp]
\centering\small
\caption{Corresponding functions across generations.}\label{tab:map}
\rowcolors{2}{acnmist}{white}
\begin{tabularx}{\linewidth}{@{}l Y Y Y@{}}
\toprule
\rowcolor{white} Function & 2G/3G & 4G (EPC) & 5G (5GC) \\
\midrule
Mobility and access control & MSC/VLR, SGSN & MME & AMF \\
Subscriber data and authentication & HLR/AuC & HSS & UDM, AUSF \\
Session control & SGSN, GGSN & MME, S-GW/P-GW control & SMF \\
User-plane anchor & GGSN & P-GW (S-GW) & UPF \\
Policy & -- (static) & PCRF & PCF \\
Radio control & BSC, RNC & eNB & gNB (CU-CP) \\
\bottomrule
\end{tabularx}
\end{table}

% ---------------------------------------------------------------------
\section{Step-by-step procedures}

\begin{procedurebox}{Analysing a new mobile system}\label{proc:analyse}
\begin{steps}
  \item Identify the subscriber database, the mobility function, the session function and the user-plane anchor.
  \item Draw the user path and count the nodes and tunnels (\cref{eq:path}).
  \item Identify the authentication method and which identities travel in clear.
  \item Identify the QoS unit (circuit, bearer, flow) and where it is enforced.
  \item Locate the interworking with the previous generation.
\end{steps}
\end{procedurebox}

\begin{procedurebox}{Planning a capacity increase}\label{proc:cap}
\begin{steps}
  \item Measure current traffic and forecast growth.
  \item Compute the required factor and split it into spectrum, efficiency and density (\cref{eq:area}).
  \item Price each lever (licences, new equipment, sites and backhaul).
  \item Choose the cheapest combination and check the cost-per-bit target (\cref{eq:cost}).
\end{steps}
\end{procedurebox}

% ---------------------------------------------------------------------
\section{Worked examples}

\begin{examplebox}{Densification needed}\label{ex:dens}
\textbf{Problem.} A hotspot must carry 10\,Tbit/s/km$^2$ with $B = 400$\,MHz and $\eta = 10$\,bit/s/Hz. How many cells per km$^2$ are needed?

\textbf{Solution.} $n = 10^{13}/(4\times10^{8}\times10) = 2500$ cells/km$^2$, one cell per 400\,m$^2$ -- an inter-site distance of about 20\,m, i.e.\ indoor
small cells.
\end{examplebox}

\begin{examplebox}{Extrapolating the trend}\label{ex:extra}
\textbf{Problem.} If the fit of \cref{sec:peak} continued, when would peak rates reach 1\,Tbit/s, a figure mentioned in 6G visions \cite{saad2020vision}?

\textbf{Solution.} $\ln(10^{12}/2\times10^{10})/0.539 = 7.3$ years after 2018, i.e.\ around 2025. That has not happened, showing that the trend is not a
law: further growth needs wider bandwidths that exist only above 100\,GHz (Lecture~26).
\end{examplebox}

\begin{examplebox}{Cost pressure}\label{ex:cost}
\textbf{Problem.} Traffic grows 25\,\% per year and revenue 2\,\%. By how much must cost per bit fall each year, and when does it halve?

\textbf{Solution.} $1.02/1.25 = 0.816$: an 18.4\,\% fall per year; $T_{1/2} = \ln 2/(\ln 1.25 - \ln 1.02) = 3.41$ years.
\end{examplebox}

% ---------------------------------------------------------------------
\section{Technology comparison}
\begin{table}[htbp]
\centering\small
\caption{Generations at a glance.}\label{tab:cmp}
\rowcolors{2}{acnmist}{white}
\begin{tabularx}{\linewidth}{@{}l Y Y Y Y@{}}
\toprule
\rowcolor{white} & 2G (GSM) & 3G (UMTS) & 4G (LTE) & 5G (NR) \\
\midrule
Multiple access & TDMA/FDMA & WCDMA & OFDMA / SC-FDMA & OFDMA, flexible numerology \\
Core & CS (+ GPRS) & CS + PS & All-IP EPC & Service-based 5GC \\
Signalling & SS7/MAP & SS7/MAP & Diameter, GTPv2-C & HTTP/2 SBI, PFCP \\
Authentication & One-way & Mutual AKA & EPS-AKA & 5G-AKA, EAP-AKA$'$, SUCI \\
QoS unit & Circuit & PDP context & EPS bearer & QoS flow, slice \\
Voice & CS & CS & VoLTE (IMS) & VoNR (IMS) \\
\bottomrule
\end{tabularx}
\end{table}

\paragraph{Same ideas, better tools} Every generation keeps the same functional roles but rebuilds them with more general technology: from
dedicated switches to IP routers to cloud software.

% ---------------------------------------------------------------------
\section{Key takeaways}
\begin{takeawaybox}
\begin{itemize}
  \item Functions persist across generations; their names, protocols and placement change.
  \item The user path flattened from four nodes (GPRS/UMTS) to one UPF in 5G.
  \item Capacity $C_A = nB\eta$: growth comes from spectrum, efficiency and densification together.
  \item Peak rate grew about 1.7 times per year from 1991 to 2018, but trends are not laws.
  \item Cost per bit must halve every $\ln 2/(\ln(1+g) - \ln(1+r))$ years -- the driver of softwarisation.
\end{itemize}
\end{takeawaybox}

% ---------------------------------------------------------------------
\section{Exercises}

\begin{exercisebox}{1}{Mapping}\label{exr:map}
Name the 4G and 5G counterparts of: (a) VLR; (b) HLR; (c) GGSN; (d) RNC.
\end{exercisebox}

\begin{exercisebox}{1}{Why packets?}\label{exr:why}
Give two reasons why mobile networks moved from circuit to packet switching for all traffic.
\end{exercisebox}

\begin{exercisebox}{2}{Growth rate}\label{exr:cagr}
LTE Rel-8 offered 300\,Mbit/s (2008) and LTE-Advanced 3\,Gbit/s (2011). Compute $k$ and the doubling time for this interval and compare with \cref{sec:peak}.
\end{exercisebox}

\begin{exercisebox}{2}{Capacity levers}\label{exr:cap}
An operator needs 20 times more area capacity. Spectrum can grow by 2.5 and efficiency by 2. What densification is needed?
\end{exercisebox}

\begin{exercisebox}{3}{Path latency}\label{exr:path}
Compare with \cref{eq:path}: (a) UMTS path with NodeB, RNC, SGSN and GGSN (1\,ms each, 300\,km fibre); (b) 5G with gNB and a UPF 30\,km away (0.3\,ms each).
\end{exercisebox}

\begin{exercisebox}{3}{The next bottleneck}\label{exr:next}
Using the five dimensions of evolution, argue which bottleneck of 5G a sixth generation is most likely to address, with a quantitative argument.
\end{exercisebox}

\subsection*{Solutions}
\addcontentsline{toc}{subsection}{Solutions}

\begin{solutionbox}{\ref{exr:map}}
(a) MME / AMF; (b) HSS / UDM (with AUSF); (c) P-GW / UPF with SMF; (d) none as a separate node -- eNB / gNB (CU-CP).
\end{solutionbox}

\begin{solutionbox}{\ref{exr:why}}
Statistical multiplexing uses capacity far better for bursty traffic, and a single IP network for voice and data lowers cost and allows Internet services
to be offered directly.
\end{solutionbox}

\begin{solutionbox}{\ref{exr:cagr}}
$k = \ln 10/3 = 0.768$ per year, $T_2 = 0.90$ years -- faster than the long-term 1.29 years, because carrier aggregation and MIMO added bandwidth and layers
in a single step.
\end{solutionbox}

\begin{solutionbox}{\ref{exr:cap}}
$20/(2.5\times2) = 4$: four times as many cells per unit area.
\end{solutionbox}

\begin{solutionbox}{\ref{exr:path}}
(a) $4\times1 + 300\,\text{km}/(2\times10^{5}\,\text{km/s}) = 4 + 1.5 = 5.5$\,ms. (b) $2\times0.3 + 0.15 = 0.75$\,ms.
\end{solutionbox}

\begin{solutionbox}{\ref{exr:next}}
Example argument: the peak-rate trend stalled because bandwidth below 100\,GHz is exhausted; $C = B\log_2(1+\mathrm{SNR})$ grows linearly in $B$ but only
logarithmically in SNR, so the next tenfold step needs tens of gigahertz of spectrum (sub-THz), plus new ways to reduce cost per bit (AI-driven
automation) and to cover the unconnected (non-terrestrial networks).
\end{solutionbox}

% ---------------------------------------------------------------------
\section{Further reading}
ITU-R M.2083 \cite{itu2015m2083} and M.2410 \cite{itu2017m2410} define the IMT-2020 vision and requirements; TS~23.002 \cite{ts23002}, TS~23.401
\cite{ts23401} and TS~23.501 \cite{ts23501} describe the architectures of the successive generations. Eberspächer et al.\ \cite{eberspaecher2009gsm},
Holma and Toskala \cite{holma2010wcdma}, Sesia et al.\ \cite{sesia2011lte} and Dahlman et al.\ \cite{dahlman2020nr} cover GSM, UMTS, LTE and NR;
Rommer et al.\ \cite{rommer2020core} cover the 5G core.

\printbibliography[heading=subbibliography,title={\ACNstrRefs}]
\end{document}
